\section{D. The boundary of the parameter disc}

We use the standard twists throughout this section. Put
\(H_*=(3/\pi)\log6\), fix a boundary parameter
\(c_*=x+iH_*\), and set \(a=x+3/2\). The distinction between a
single degenerate fibre and the total quotient over its root disc is
essential below. All topological group quotients have their quotient
topology; a Hausdorff reflection means a continuous quotient universal
for continuous maps to Hausdorff spaces.

\subsection{D.1. The intrinsic obstruction to extension}

\begin{proposition}[D.1.1; Proved]
For \(0<|w|<216\), the reduced multiplicity-four surface has
\[
 j(\operatorname{Alb}S_2(w))=J(iw/216),\qquad
 J(q)=q^{-1}+744+196884q+\cdots.
\]
This holomorphic function has no meromorphic continuation through any
point of \(|w|=216\).
\end{proposition}
\begin{proof}
The fixed-coordinate periods on the product cover, including the
isogeny kernel and the affine deck translation, generate
\(\frac12\mathbb Z+\frac{\kappa_2}{4}\mathbb Z\).
Multiplication by two therefore identifies the Albanese lattice with
\(\mathbb Z+(\kappa_2/2)\mathbb Z\). This verifies the precise step of
R7 used here; taking only the invariant eigenspace would omit the
half-period. The values in R3 and \(\mu(z_2)=(1-i)/2\) give
\[
 \kappa_2=c+\frac32-iH_*,\qquad
 e^{-\pi i\kappa_2}=e^{-3\pi i/2}e^{-\pi H_*}e^{-\pi ic}
 =\frac{iw}{216}.
\]
The positive-imaginary representative is \(-\kappa_2/2\), proving the
identity with the stated Fourier normalization.

For completeness we reproduce the obstruction in the earlier return,
C.3.4, rather than regard a numerical Fourier plot as evidence. At a
reduced rational \(r=b/d\) choose a matrix in
\(\operatorname{SL}_2(\mathbb Z)\) with bottom row \((d,-b)\).
For \(z=r+iy\), its transformed imaginary part is \(1/(d^2y)\).
Modular invariance and the convergent Fourier series imply
\(\log|j(r+iy)|=2\pi/(d^2y)+o(1)\).
A meromorphic function near \(r\) grows at most polynomially in
\(1/y\), so no continuation exists there. Every real interval contains
such a rational, excluding continuation at irrational points as well.
A local logarithm of \(q=iw/216\) transfers this result to the circle.
This is a derivation of the earlier natural-boundary result, with the
normalizing constant checked exactly.
\end{proof}

\begin{lemma}[D.1.2; Proved]
Let \(V\) be a connected Riemann surface and
\(\pi:\mathcal X\to V\) a flat proper holomorphic map of reduced
complex spaces. Suppose that over a nonempty open subset \(V_0\), the
fibres are smooth compact threefolds whose unique anticanonical divisor
is \(2S_v\), where \(S_v\) is a smooth connected surface with
\(b_1(S_v)=2\). Then there is an open subset \(V^{\rm good}\subset V\), with at
most countable complement, on which the reduced support of the
relative anticanonical divisor is a smooth proper surface family. Its
Albanese elliptic invariant is holomorphic.
\end{lemma}
\begin{proof}
Flatness and the smoothness of the fibres over \(V_0\) imply that
\(\pi\) is a submersion there: the local analytic smoothness criterion
applies because the base is smooth. The relative nonsmooth locus is a
closed analytic subset, by the Jacobian criterion. Its proper image is
analytic by Grauert (1960, \S7, Satz~1), the proper mapping theorem.
It misses \(V_0\), hence is discrete. We remove it and work over the
connected remaining curve \(V'\). Removing a discrete subset from a
connected Riemann surface does not disconnect it. Thus
\(\mathcal X'\to V'\) is a smooth proper map.

Write \(L=K_{\mathcal X'/V'}^{-1}\). Grauert's direct-image theorem
(1960, \S6, Hauptsatz~I) makes \(F=\pi_*L\) and \(R^1\pi_*L\)
coherent. We prove the needed base change directly over the smooth
curve; we do not use the originally stated \S7, Satz~2, whose claimed
closed exceptional locus was corrected in Grauert's 1963 erratum.
For a local base parameter \(t\) vanishing at \(v\), smoothness and
invertibility of \(L\) give an exact sequence
\(0\to L\xrightarrow{t}L\to i_*L_v\to0\), where
\(i:X_v\hookrightarrow\mathcal X'\) is the closed embedding.
Its direct-image sequence gives, at \(v\),
\[
 0\longrightarrow F_v/tF_v\longrightarrow H^0(X_v,L_v)
 \longrightarrow (R^1\pi_*L)_v[t]\longrightarrow0.
\]
The last term denotes elements annihilated by \(t\). The support of
the torsion subsheaf of the coherent module \(R^1\pi_*L\) is closed
discrete, by the structure theorem for finitely generated modules
over the local discrete valuation rings of the smooth base curve.
Outside that support the canonical base-change map is an isomorphism.
Also \(F\) is torsion-free: multiplication by a nonzero base function
is injective on sections of \(L\). It is therefore locally free.
Its rank is one by choosing an interior point outside the discrete
torsion support. These statements concern the canonical map, not
merely equal dimensions.
Locally a frame gives a relative section \(\sigma\) of \(L\),
nowhere identically zero on a fibre. Its divisor is independent of the
frame. Let \(D\) be its reduced support. On \(V_0\), the divisor is
\(2S_v\) and \(D\) is exactly \(S_v\).

The map \(D\to V'\) is proper, since \(D\) is closed in a proper
space. There is no vertical divisor component: a vertical divisor in
the smooth total space would contain a whole irreducible fibre,
contrary to the nonzero restriction of \(\sigma\). Every horizontal
component has proper image all of the connected base. Since the
interior fibre is irreducible, there is exactly one horizontal
component, and its multiplicity in \(\operatorname{div}\sigma\) is
two. Thus \(\operatorname{div}\sigma=2D\) on this base. In the smooth
ambient space \(D\) is Cartier. It is flat over the curve because
its local rings are torsion-free over the base discrete valuation
ring. Restriction of the Cartier identity to an interior fibre gives
\(2D_v=2S_v\), whence \(D_v=S_v\) with its reduced structure.
Flatness and the smoothness of these fibres make \(D\to V'\) smooth
over \(V_0\). Its analytic nonsmooth locus therefore misses \(V_0\);
the proper mapping theorem once more makes its image discrete.
After deleting that image, \(D\to V^{\rm good}\) is a smooth proper
family of connected surfaces. Each deletion is closed discrete in the
then-current open curve. Their finite union is at most countable,
and the final good locus is open; we do not assert that successive
exceptional sets cannot accumulate at an earlier removed point.

Ehresmann's theorem identifies the smooth fibres along every path;
in particular their first Betti number is two. Buchdahl (1999,
Theorem~11) then gives a K\"ahler metric on every fibre. Hence
\(h^{1,0}=1\), every holomorphic one-form is closed, and the integral
period map has a rank-two lattice image. Constant \(h^{1,0}\),
properness, and analytic base change supply a local relative
holomorphic one-form. Locally on a simply connected
base disc, a relative one-form and a basis of the free part of
\(H_1\) give two holomorphic periods. Their ratio lies in one
half-plane after orienting the basis. To see holomorphicity directly,
choose local primitives of the closed relative form on simply
connected fibre charts; their differences are holomorphic functions
of the base, and the finite sum around a transported cycle is its
period. The modular invariant of the ratio is independent of all
choices and is holomorphic. This proves the last assertion without
requiring a global choice of section of the Albanese family.
\end{proof}

\begin{theorem}[D.1.3; Proved; Lead 8]
For every \(w_0\) with \(|w_0|=216\) and every connected open
neighbourhood \(U\) of \(w_0\), there is no flat proper holomorphic
map \(\mathcal X\to U\) from a reduced complex space whose fibres
over \(U\cap\{0<|w|<216\}\) are biholomorphic, fibre by fibre, to
\(X_w\). In particular an extension isomorphic to the given family
over that intersection does not exist.
\end{theorem}
\begin{proof}
We may restrict a putative extension to a small disc meeting the
circle, with closure away from zero. This is a holomorphic base
restriction; properness and flatness are preserved. The calculation
of the canonical evaluation in the earlier return, A.1.1, gives
\(K_{X_w}=\mathcal O(-2S_2(w))\) and
\(h^0(K_{X_w}^{-1})=1\). Its sole divisor is \(2S_2(w)\), so it is
intrinsic under every fibrewise biholomorphism. The surface has
\(b_1=2\): on its finite product cover the rotation kills the moving
factor in invariant first cohomology and fixes the elliptic factor.
These are the exact R1 and R7 steps required in Lemma~D.1.2.

The at most countable exceptional set in that lemma cannot contain
an arc of the circle. Choose a boundary point outside it and then a smaller disc
on which the reduced divisor is a smooth proper surface family. Its
holomorphic Albanese invariant agrees, pointwise on the inner part,
with \(J(iw/216)\). It is therefore an analytic continuation through
that point, contrary to Proposition~D.1.1. No choice of the
fibrewise biholomorphisms was used. Thus the weaker hypothesis already
leads to the contradiction.
\end{proof}

\subsection{D.2. The local quotient and its Hausdorff reflection}

\begin{lemma}[D.2.1; Proved]
Near the order-four root centre there is a holomorphic coordinate
\[
 h=\mu-(1-\tau)/2
\]
with nonzero derivative in \(s_2\). In the fixed complex-linear fibre
coordinates \(\eta=\zeta_1\), \(\xi=\zeta_2+\zeta_1/2\), and the
integral ordered basis
\((\widehat w,\widehat u,\widehat\delta,
\widehat\gamma+3\widehat u-3\widehat w)\), the periods at \(c_*\) are
\begin{equation}\label{eq:d-period-normal}
 (1,\tfrac12),\quad(\tau,\tfrac12+h),\quad(0,1),\quad(6h,a+B(h)),
\end{equation}
where
\[
 \tau=i-\frac{3\pi i}{2}h^2+O(h^4),\qquad
 B(h)=-3ih^2+O(h^4),\qquad
 D_{c_*}=-3|h|^2+O(|h|^4).
\]
For \(c\) near \(c_*\), only the last \(\xi\)-period changes, by
\(c-c_*\).
\end{lemma}
\begin{proof}
The basis change has determinant of absolute value one. Substitution
into the original period matrix gives \eqref{eq:d-period-normal}
exactly, with
\(B=\beta+6h+3(1-\tau)/2-a\).
At the centre \(B(0)=0\). Put
\(G(h,\tau)=\wp_\tau((1-\tau)/2+h)\) and
\(H(h,\tau)=\log(\vartheta_3(h,\tau)/\eta_D(\tau))\).
The half-period identities and theta heat equation give
\[
 G(0,i)=0,\quad G_{h^2}(0,i)=-g_2(i)/4,\quad
 G_\tau(0,i)=\frac{i g_2(i)}{6\pi},
\]
\[
 H(0,i)=\log\sqrt2,\qquad H_\tau(0,i)=0,\qquad
 H_{h^2}(0,i)=-\pi/2.
\]
Here \(G_{h^2}\) is the power-series coefficient. For the derivative
constant, write \(A=\vartheta_2(0,i)^4=\vartheta_4(0,i)^4\).
Then \(g_2(i)=4\pi^4A^2\), while differentiating
\(e_2=(\pi^2/3)(\vartheta_2^4-\vartheta_4^4)\) with the heat
identities gives \(e_2'(i)=2\pi^3iA^2/3\).
The logarithmic derivatives used in this differentiation are
\[
 \frac{\partial_\tau\vartheta_2}{\vartheta_2}
  =\frac{\pi i}{12}(E_2+\vartheta_3^4+\vartheta_4^4),\qquad
 \frac{\partial_\tau\vartheta_4}{\vartheta_4}
  =\frac{\pi i}{12}(E_2-\vartheta_2^4-\vartheta_3^4).
\]
They follow by applying the theta heat equation to the half-period
identities. The coefficient arithmetic is independently checked in
\texttt{d\_push\_exact.py}. In particular \(G_\tau(0,i)\ne0\).
The holomorphic implicit-function theorem and evenness of \(G\) give
\(\tau=\tau(h)\), even in \(h\), with the displayed quadratic
coefficient. Since \(\tau-i\) has order two in the root coordinate,
\(h\) has order one.

The identity verified in the return's C.2.2 is
\[
 \log\mathcal N=\operatorname{Re}H-
       \frac{\pi(\operatorname{Im}h)^2}{\operatorname{Im}\tau}.
\]
The displayed coefficients, \(\tau-i=O(h^2)\), and evenness give
\(\log\mathcal N-\log\sqrt2=-(\pi/2)|h|^2+O(|h|^4)\).
Multiplication by \(6/\pi\), using the boundary normalization, gives
the formula for \(D_{c_*}\). Directly from the definition,
\[
 D_{c_*}=\operatorname{Im}B-
             6(\operatorname{Im}h)^2/\operatorname{Im}\tau.
\]
Consequently \(\operatorname{Im}B=-3\operatorname{Re}(h^2)
+O(|h|^4)\). Since \(B\) is holomorphic and \(B(0)=0\), its Taylor
coefficients force \(B=-3ih^2+O(h^4)\). This argument checks the
quadratic coefficient that the first-order Lead~9 omitted.
\end{proof}

\begin{proposition}[D.2.2; Refuted; Lead 9]
The assertion in Lead~9 that total-space nonseparation is confined to
\(\mathbb C(2,-1)^t\), modulo central periods, is false. At the
boundary value \(a=0\), the central points with \((\eta,\xi)=(0,0)\)
and \((0,-3i)\) cannot be separated in the root-family quotient.
They are not related by that complex line and the central periods.
\end{proposition}
\begin{proof}
Choose \(h_n=1/n\), which defines actual root points by Lemma~D.2.1,
and take the integral lattice vector
\[
 \lambda_n=n^2(\widehat\gamma+3\widehat u-3\widehat w)
            -6n\widehat w+3n\widehat\delta.
\]
Its period is exactly \((0,n^2B(1/n))\), and this tends to
\((0,-3i)\). Thus the orbit of \((h_n,0)\) also has a representative
converging to \((0,0,-3i)\). The two central quotient points are
nonseparated. In these coordinates the asserted complex line is
\(\{\xi=0\}\), and every central period has real \(\xi\)-coordinate;
therefore their difference is not in the asserted relation. The
quadratic model calculation and its integer coefficients are checked
exactly in \texttt{d\_push\_exact.py}; the analytic error is
\(n^2O(n^{-4})\to0\), so the counterexample is for the actual family.
\end{proof}

\begin{theorem}[D.2.3; Proved]
Let \(\mathcal T=(\Delta_h\times\mathbb C^2)/\Lambda\) at the fixed
boundary parameter \(c_*\), after shrinking \(\Delta_h\). Two
central points represented by \((\eta,\xi)\) and
\((\eta',\xi')\) are nonseparated precisely when
\[
 \operatorname{Re}(\xi'-\xi)\in G_a
       :=\tfrac12\mathbb Z+a\mathbb Z\quad(a\in\mathbb Q),
\]
and every pair is nonseparated when \(a\notin\mathbb Q\).
There is no nonseparation between distinct base fibres, or between
distinct points of a noncentral fibre. The Hausdorff reflection
\(\mathcal T^{\mathrm H}\to\Delta_h\) has the original complex tori
over \(h\ne0\), and its central fibre is
\[
 \mathbb R/G_a\simeq S^1\quad(a\in\mathbb Q),\qquad
 \{*\}\quad(a\notin\mathbb Q).
\]
If \(a=m/n\) is in lowest terms, \(n>0\), then
\(G_a=\operatorname{lcm}(2,n)^{-1}\mathbb Z\).
\end{theorem}
\begin{proof}
We give the full limiting-group argument in D.3.1; its formula also
applies to the fixed boundary family with \(c=c_*\), \(h\ne0\).
Put \(d_0=\operatorname{lcm}(2,n)\) in the rational case. In the dual
lattice, the fixed integral covector whose central limit is
\(d_0\operatorname{Re}\xi\) has
\((p,q,r)=(0,0,d_0)\). The exact dual formula
\eqref{eq:d-dual-exact} below shows that it extends continuously, as a
real-linear functional \(\ell_h\) on \(\mathbb C^2\), to that limit:
its last coefficient is \(O(h^2)\), its other moving coefficients
are \(O(h)\). Therefore
\[
 (h,z)\longmapsto\exp(2\pi i\ell_h(z))
\]
is continuous, invariant under every period, and restricts at the
centre to \(\exp(2\pi i d_0\operatorname{Re}\xi)\).
Points with different indicated central classes are separated by this
map to the circle.

For the converse take a ray \(h=\epsilon\omega\),
\(\epsilon\downarrow0\), whose real and imaginary parts have
irrational ratio. More precisely choose \(\omega\) so that
\(p\operatorname{Re}\omega+q\operatorname{Im}\omega\ne0\)
for every nonzero pair \((p,q)\) in the countable half-integral
lattice in D.3.1. The asymptotic dual formula shows that all such
covectors escape to infinity. For rational \(a\), the surviving dual
group is exactly \(\mathbb Z d_0\operatorname{Re}\xi\); for
irrational \(a\), it is zero. Pontryagin--Chabauty duality
(Cornulier, Theorem~3.1) now gives the period-group limits
\[
 \{z:d_0\operatorname{Re}\xi(z)\in\mathbb Z\},\qquad
 \mathbb C^2,
\]
respectively. Every element of the indicated limit is the limit of
actual period vectors. Such a sequence supplies intersecting images
of every two prescribed central neighbourhoods, proving
nonseparation. This includes the quadratic-scale differences of
D.2.2, not only first-order differences.

Collapse exactly these central classes and leave every other fibre
unchanged. The resulting quotient is Hausdorff. The base coordinate
separates different fibres. Within a noncentral fibre, continuous
functions supported over a compact subdisc away from zero separate
distinct points, because there the original action is properly
discontinuous and gives a smooth torus bundle. In the rational central
fibre the displayed circle-valued function separates different
classes. These functions descend through the proposed quotient, so
it is Hausdorff. Every continuous map from \(\mathcal T\) to a
Hausdorff space identifies nonseparated pairs and therefore factors
uniquely and continuously through this quotient. This proves the
universal property, including its topology. The elementary greatest
common divisor calculation gives the stated value of \(G_a\).
\end{proof}

\begin{proposition}[D.2.4; Proved]
At the boundary, the construction over \(B\setminus\{1\}\) is a
complex manifold, since R3 makes every remaining finite period map
invertible and the cusp completion is unchanged.
The order-four deck action descends continuously to
\(\mathcal T^{\mathrm H}\). On its rational central circle it is
translation by \(-a/4\); on the irrational central point it is
trivial. Its finite quotient is Hausdorff with central fibre a circle
or a point, respectively. Neither quotient is a complex space over
the base compatible with the given holomorphic family off the centre.
The linear foliation by \(\mathbb C(2,-1)^t\) on the original central
torus cover has compact elliptic leaves. For rational \(a\) the
leaves are closed; for irrational \(a\) the closure of a leaf is the
set with its fixed value of \(\operatorname{Im}\xi\).
\end{proposition}
\begin{proof}
The deck transformation is a biholomorphism of the original quotient,
so the universal property extends it uniquely to a homeomorphism of
the Hausdorff reflection. The action on the fixed product coordinate
at the centre is \(\xi\mapsto\xi-\kappa_2/4=\xi-a/4\).
A finite group quotient of a Hausdorff space is Hausdorff: finite
orbits are separated by intersecting finitely many translated open
neighbourhoods. A finite rotation quotient of a circle is a circle.

We verify properness of the reflected map, rather than assume it.
On a closed small base disc every noncentral torus point has a
representative in the image of \([0,1]^4\) under its period matrix.
These representatives belong to a single compact subset of
\(\mathbb C^2\). At the centre their images meet every class of the
circle or the point. The image of that compact parameterized
parallelepiped is therefore the full inverse image of the closed
disc in the reflection. The reflected map is proper; the same holds
after the finite quotient. If it were holomorphic on a compatible
complex space, the central fibre would be analytic. A circle cannot
be a complex analytic space: a positive-dimensional analytic germ
has real dimension at least two, and a compact zero-dimensional
analytic space is finite. A central point is also impossible here:
upper semicontinuity of fibre dimension, or the local dimension
inequality at a point in the closure of the three-dimensional smooth
locus, forces a fibre dimension of at least two for a holomorphic
map to a curve. Thus this Hausdorff reflection cannot be glued as a
complex-space filling to the remaining constructed manifold.

For the leaf assertion, the moving line is \(\{\xi=0\}\).
At the centre the projection of the period group to the \(\xi\)-line
is \(G_a\subset\mathbb R\). Its intersection with the moving line
is a rank-two lattice in \(\mathbb C_\eta\): the periods are
\(p+qi\) with \((p+q)/2+ka+\ell=0\) for integers
\(p,q,k,\ell\). For irrational \(a\) this is the index-two lattice
\(p+q\equiv0\pmod2\); for rational \(a\) it is either that lattice
or \(\mathbb Z[i]\). Both give intrinsic elliptic leaf manifolds isomorphic to
\(E_i\). In the irrational case their maps into the non-Hausdorff
central quotient need not be topological embeddings; ``leaf'' here
means the image with this intrinsic immersed topology. When \(G_a\)
is discrete, the saturated leaf is closed.
When it is dense in \(\mathbb R\), its saturation has closure
\(\mathbb C_\eta\times(\operatorname{Re}\xi+\mathbb R)\), with
\(\operatorname{Im}\xi\) fixed. This is a statement in the original,
possibly non-Hausdorff central fibre. Its own Hausdorffization is the
rank-three quotient described in R9; it is not the central fibre of
\(\mathcal T^{\mathrm H}\).
\end{proof}

\subsection{D.3. All local limiting period groups}

\begin{theorem}[D.3.1; Proved]
Let \(c_n\to c_*\) with \(\operatorname{Im}c_n<H_*\), and let
\(h_n\to0\). Set
\[
 \delta_n=-D_{c_n}(h_n)>0,\qquad
 V_n=\left(\frac{6\operatorname{Re}h_n}{\delta_n},
           \frac{6\operatorname{Im}h_n}{\delta_n},
           \frac{\operatorname{Re}(c_n-c_*)}{\delta_n}\right).
\]
Define the discrete group
\[
 L_a=\{(p,q,r)\in\mathbb R^3:
    r\in\mathbb Z,\ p+r/2\in\mathbb Z,
    \ q+r/2\in\mathbb Z,\ ar\in\mathbb Z\}.
\]
It has rank three for rational \(a\), and is
\(\mathbb Z^2\times\{0\}\) for irrational \(a\).
Every convergent subsequence of the closed lattices
\(\Gamma_n=\Pi_{c_n}(h_n)\mathbb Z^4\subset\mathbb C^2\) has the
following form. There is a primitive subgroup \(K\subset L_a\) and
a homomorphism \(\ell:K\to\mathbb R\) such that
\begin{equation}\label{eq:d-limit-group}
 \Gamma_\infty=
 \{(\eta,\xi):p\operatorname{Re}\eta+q\operatorname{Im}\eta
       +r\operatorname{Re}\xi+\ell(p,q,r)\operatorname{Im}\xi
       \in\mathbb Z\text{ for every }(p,q,r)\in K\}.
\end{equation}
Here
\[
 K=\{\lambda\in L_a:\lambda\cdot V_n
                      \text{ has a finite limit}\},\qquad
 \ell(\lambda)=\lim_n\lambda\cdot V_n,
\]
and for \(\lambda\notin K\) the absolute value tends to infinity.
Conversely every primitive \(K\subset L_a\) and every such
\(\ell\) occur for sequences with \(\operatorname{Im}c_n<H_*\).
Thus the possible pointed additive quotient limits are real tori of
any dimension from zero to three for rational \(a\), and from zero
to two for irrational \(a\), with the marked groups given explicitly
by \eqref{eq:d-limit-group}.
\end{theorem}
\begin{proof}
We use the fixed complex-linear coordinates of D.2.1. Write
\(h=X+iY\), \(\tau=R+iT\), and
\(z=B(h)+c-c_*\). The real period matrix, with the basis of D.2.1,
is
\[
 A=\begin{pmatrix}
 1&R&0&6X\\
 0&T&0&6Y\\
 1/2&1/2+X&1&a+\operatorname{Re}z\\
 0&Y&0&\operatorname{Im}z
 \end{pmatrix},\qquad
 \det A=T D_c.
\]
For an integral dual label \((m_w,m_u,m_d,m_g)\), put
\(r=m_d\), \(p=m_w-r/2\), \(q=m_u-r/2\).
The real-linear covector \(A^{-t}(m_w,m_u,m_d,m_g)^t\) is
\begin{equation}\label{eq:d-dual-exact}
 \left(p,\frac{q-pR-rX-tY}{T},r,t\right),\quad
 t=\frac{m_g-ra-6pX-6Y(q-pR-rX)/T-r\operatorname{Re}z}{D_c}.
\end{equation}
Multiplication by \(A^t\) proves the formula; the entire equation is
checked symbolically in \texttt{d\_push\_exact.py}.

A bounded sequence of dual vectors has bounded integral labels,
since \(A_n\) is bounded. After subselection its label is constant.
Its limit must satisfy \(m_g=ra\), because
\(A_0^t\) has that relation. The remaining labels are exactly
\(L_a\). Conversely consider a fixed label from \(L_a\).
Lemma~D.2.1 gives
\[
 T=1+O(|h|^2),\quad R=O(|h|^2),\quad
 \operatorname{Re}B=6XY+O(|h|^4),
\]
\[
 \delta_n=-\operatorname{Im}(c_n-c_*)
                       +3|h_n|^2+O(|h_n|^4).
\]
In particular \(\delta_n\ge C|h_n|^2\) for a fixed \(C>0\)
after shrinking the disc. The numerator in
\eqref{eq:d-dual-exact}, with \(m_g=ra\), is
\[
 -6pX-6qY-r\operatorname{Re}(c-c_*)
       +O(|h|^3)(|p|+|q|+|r|).
\]
The cancellation of the quadratic real term is essential. Division
by \(D_c=-\delta_n\) makes the error \(O(|h_n|)\). Thus
\[
 t_n=(p,q,r)\cdot V_n+o(1).
\]
Whenever this is bounded, the other three coefficients converge to
\((p,q,r)\). This proves that the dual limit, if it exists, is the
discrete graph
\(\{(p,q,r,\ell(p,q,r)):(p,q,r)\in K\}\).
There can be no further elements from varying labels, because a
bounded sequence of labels lies in a finite set. A convergent dual
group cannot contain two different limits with the same label: their
difference would have label zero, whose sole preimage under every
\(A_n^t\) is zero. Nor can a surviving label escape along a
subsequence, since an element of a Chabauty limit must be approximable
along the whole convergent sequence. Consequently every label either
has the indicated finite limit or escapes to infinity. Addition and
integer division show that \(K\) is a subgroup and is primitive.

Cornulier's Theorem~3.1 identifies convergence of closed subgroups of
a finite-dimensional real vector space with convergence of their
Pontryagin annihilators. Applying it to the dual graph gives exactly
\eqref{eq:d-limit-group}. This theorem's hypotheses are met by the
real vector space underlying \(\mathbb C^2\) and by the closed full
lattices \(\Gamma_n\); no nonclosed group is used as a Chabauty
point.

To prove realizability, extend \(\ell\) linearly to a covector
represented by a vector \(v_0\in\mathbb R^3\). Since \(K\) is
primitive, \(L_a\cap\operatorname{span}_{\mathbb R}K=K\).
Choose \(v\) orthogonal to that span and not orthogonal to any member
of \(L_a\setminus K\); the excluded proper hyperplanes are
countably many, and cannot cover the orthogonal space. If
\(K=L_a\), take \(v=0\). Put \(V_n=nv+v_0=(A_n',B_n',C_n')\).
Choose positive \(d_n\to0\) with
\(d_n(1+\|V_n\|^2)\to0\), and define
\[
 h_n=d_n(A_n'+iB_n')/6,\quad x_n=d_nC_n',\quad
 \epsilon_n=d_n+D_{c_*}(h_n)>0,
 \qquad c_n=c_*+x_n-i\epsilon_n.
\]
The inequality follows from \(D_{c_*}(h)=-3|h|^2+O(|h|^4)\).
These choices give \(-D_{c_n}(h_n)=d_n\) exactly, so their vectors
are the prescribed \(V_n\). On \(K\) the scalar products converge
to \(\ell\), and off \(K\) they escape. This constructs every
claimed limit inside the admissible parameter domain.

Finally the dual graph is a discrete free group of rank
\(\operatorname{rank}K\). Its annihilator quotient is a compact
connected real torus of that dimension. In the fixed Euclidean metric,
\(d_n([z],[z'])=\operatorname{dist}(z-z',\Gamma_n)\).
Chabauty convergence gives local uniform convergence of these
functions: minimizers relevant on a bounded set can be restricted to
a bounded set, and the two defining Chabauty approximation properties
then give both inequalities. This identifies the stated pointed
quotient limits. They are assertions about real topological and
metric groups, not complex tori of a varying complex dimension.
\end{proof}

\begin{corollary}[D.3.2; Proved]
There is no path-independent limiting group at the boundary even
when \(c_n\to c_*\) from the admissible side. At \(h_n=0\) and
\(c_n=c_*-i\epsilon_n\), all of \(L_a\) survives, giving a
three-dimensional real torus if \(a\) is rational and a
two-dimensional real torus if \(a\) is irrational. Other admissible
paths give a point. The fixed-boundary Hausdorff reflection of D.2.3
therefore cannot be identified with all pointed fibre limits.
\end{corollary}
\begin{proof}
The first sequence has \(V_n=0\), hence \(K=L_a\) and \(\ell=0\).
For the second, take \(K=0\) in the explicit realization in D.3.1.
The limits have different real dimensions. Away from \(t=1\), the
periods stay invertible and the local complex families converge to
their ordinary value at \(c_*\); the ambiguity is at the specified
root centre. No statement of convergence of compact complex
threefolds follows from a real-group limit. In particular D.1.3
already excludes the proposed holomorphic proper extension.
\end{proof}

\subsection{D.4. The supercritical locus}

\begin{proposition}[D.4.1; Proved]
For \(\operatorname{Im}c=H_*+\epsilon\), with
\(0<\epsilon\ll1\), the zero locus \(D_c=0\) is one real-analytic
simple closed curve surrounding \(t=1\). In the coordinate of D.2.1,
\[
 |h|^2=\epsilon/3+O(\epsilon^2).
\]
In the fixed coarse coordinate it has the leading shape
\[
 |t-1|=\frac{\pi^4E_4(i)}4\epsilon^2
                         \bigl(1+O(\epsilon)\bigr).
\]
The error is uniform along the curve. The region containing \(t=1\)
has \(D_c>0\); the region containing the cusp has \(D_c<0\).
\end{proposition}
\begin{proof}
The Hessian of \(D_{c_*}(h)=-3|h|^2+O(|h|^4)\) is negative
at zero. Its radial derivative is \(-6|h|+O(|h|^3)\), uniformly
in angle. Consequently each sufficiently small positive level of
\(-D_{c_*}\) has exactly one radius in each direction and is a
real-analytic simple closed curve, with the displayed expansion.
The order-four deck action preserves it, acts freely on it, and the
coarse map \(t\) is a four-fold covering of its image. Since the
punctured root quotient is a punctured disc, that image is again a
simple closed curve. The strict global maximum in R3, including the
independent certificate in Part E.4, gives a uniform gap on the
complement of any fixed root neighbourhood; the cusp tends to minus
infinity. Thus no other zero-locus component exists for small enough
\(\epsilon\).

Ramanujan's differential identities give
\(j''(i)=-2\pi^2j(i)E_4(i)\): indeed
\(j'=-2\pi i jE_6/E_4\) and
\(E_6'(i)=-\pi iE_4(i)^2\). Since \(j=1728t\), Lemma~D.2.1 gives
\[
 t-1=\frac{9\pi^4E_4(i)}4h^4+O(h^6).
\]
Insertion of the radial expansion proves the coarse-coordinate
formula, including its order of error. The sign statement follows
from the strict radial decrease near zero and from the cusp sign.
\end{proof}

\begin{proposition}[D.4.2; Proved]
For any height \(H=\operatorname{Im}c>H_*\), the zero locus is the
well-defined modular-norm level set
\[
 Z_H=\left\{t:\mathcal N(\mu(t),\tau(t))
                    =2\sqrt3\,e^{-\pi H/6}\right\},\qquad
 \wp_{\tau(t)}(\mu(t))=0.
\]
It is compact away from the cusp, and is a real-analytic set in the
finite root charts. At every regular point it is a smooth real curve.
The construction gives a complex manifold precisely over
\(B\setminus Z_H\), retaining the original Mumford fibre at the cusp.
The open sublevel set \(\{D_c<0\}\), with the cusp adjoined, is
connected. At every point of \(Z_H\), the period quotient is the
rank-three quotient of R9; its exact Hausdorffization is described
below.
\end{proposition}
\begin{proof}
The formula for \(D\) in R3 gives the equation, and the theta
transformation law shows that it is independent of the modular and
period representatives. Compactness away from the cusp follows from
\(D_{\rm part}\to-\infty\). The analytic and regular-level
assertions follow in each root chart from the real-analytic implicit
function theorem. On the complement, the real period matrix is
invertible. On compact subsets of that complement its inverse is
bounded, which proves uniform discreteness and proper discontinuity
of the lattice action. The quotient is therefore a complex manifold
with compact torus fibres. The finite deck action is free for the
same invariant-coordinate reason as R1: a nontrivial power translates
that circle by a noninteger fraction, regardless of the sign of
\(D\). Near the cusp \(D<0\), so its existing holomorphic Mumford
completion and its overlaps remain available.
At a zero with irrational kernel line, the fibre quotient is already
non-Hausdorff. At a zero with rational kernel choose a primitive
integral kernel vector. Its holomorphic period vanishes at that root
point but is not identically zero nearby. Its first nonzero Taylor
coefficient, multiplied by arbitrarily large integer multiples and
approached along suitable root directions, yields nonseparated
central points. Thus no neighbourhood of a zero in the naive quotient
is a Hausdorff complex manifold. This justifies the asserted maximal
open base locus.

For connectedness, in any holomorphic coordinate \(z\) with
\(T=\operatorname{Im}\tau>0\), differentiation gives
\[
 \partial_z\partial_{\bar z}D_{\rm part}
 =-\frac3T\left|\mu'(z)-
         \frac{\operatorname{Im}\mu}{T}\tau'(z)\right|^2\le0.
\]
Thus the function is superharmonic on the root charts. This property
and the minimum principle descend under the finite quotient; one
can equally apply the principle on a root cover of a putative
minimum. A component of the strict sublevel set not containing the
cusp would have compact closure with boundary at level \(-H\), and
would attain a smaller interior minimum. The strong minimum
principle excludes this. A sufficiently small cusp neighbourhood is
contained in the sublevel set, so all components reaching the cusp
are the same one. Hence the sublevel set is connected.
\end{proof}

\begin{proposition}[D.4.3; Proved]
At a zero of \(D_c\), let \(A:\mathbb R^4\to\mathbb C^2\) be
the real period map and \(\ker A=\mathbb R k\), with the kernel
vector stated in R9. Define
\[
 L_k=\{m\in\mathbb Z^4:m\cdot k=0\},\qquad r=\operatorname{rank}L_k,
 \qquad V_k=\bigcap_{m\in L_k}\ker m.
\]
Then
\[
 \overline{A\mathbb Z^4}=A(\mathbb Z^4+V_k),\qquad
 \mathbb C^2/\overline{A\mathbb Z^4}\simeq
                         \mathbb R\times(S^1)^r
\]
as real topological groups. The original quotient is a complex Lie
group precisely when \(k\) is rational, in which case \(r=3\).
Otherwise it is non-Hausdorff. Every Chabauty limit of nearby full
lattices contains \(\overline{A\mathbb Z^4}\); the quotient topology
at the fixed zero and a limit of full lattices need not coincide.
\end{proposition}
\begin{proof}
The last three period columns are real-linearly independent, so the
rank is exactly three. This confirms the linear algebra of R9
without a positivity assumption. The annihilator characters of
\(A\mathbb Z^4\) are the real covectors \(\lambda\) with
\(A^t\lambda\in\mathbb Z^4\). Their integral labels belong exactly
to \(L_k\). The smallest rational subspace containing \(k\) is
\(V_k\); the closure of the image of \(\mathbb Z^4\) in
\(\mathbb R^4/\mathbb R k\) is its image after adding \(V_k\).
For an elementary verification, the two groups have the same
annihilator characters, and characters separate closed subgroups of
a real vector space by putting such a subgroup into the form
\(\mathbb R^d\times\mathbb Z^e\). This proves the closure formula.
Inside the three-dimensional image hyperplane its connected part
has dimension \(3-r\), and its discrete quotient has rank \(r\).
The one transverse real direction gives the displayed group.

The image is discrete exactly when the kernel line has a primitive
integral generator: extend such a generator to an integral basis for
the converse; for necessity, a discrete subgroup of a real
three-space has rank at most three, forcing a nonzero integral
kernel. A nondiscrete finitely generated subgroup cannot be closed,
since a nondiscrete closed subgroup contains an uncountable vector
group. Finally each fixed integral period converges to its limiting
period; a Chabauty limit is closed and therefore contains their
closure. This does not identify it with that closure.
\end{proof}

\begin{proposition}[D.4.4; Proved]
On an ordinary modulus chart where \(P=\wp'_\tau(\mu)\ne0\), put
\[
 L=\partial_\mu\log\vartheta_1(\mu,\tau),\qquad
 Z=L+2\pi i\frac{\operatorname{Im}\mu}{\operatorname{Im}\tau}.
\]
Then the exact real-analytic derivative is
\[
 \partial_\tau D_{\rm part}
       =\frac{iZ}{4\pi^2P}(3ZP-2g_2).
\]
Away also from the ramification of \(j\), the critical-point equation
on the coarse base is therefore \(Z=0\) or \(3ZP=2g_2\).
\end{proposition}
\begin{proof}
Put \(E=\pi^2E_2/3\). The theta heat equation,
\(\wp=-\partial_\mu L-E\), the differential equation
\(\wp''=6\wp^2-g_2/2\), and
\(4\pi i E'=-2E^2+g_2/6\) give, with \(\mu\) held fixed,
\[
 \partial_\tau\wp=
 \frac{2\wp^2-2E\wp+L\wp'-g_2/3}{2\pi i}.
\]
Indeed differentiating the heat equation twice in \(\mu\) gives
\((\wp''-2(\wp+E)^2+2L\wp')/(4\pi i)-E'\), which simplifies to
the displayed formula. Implicit differentiation of \(\wp(\mu)=0\)
therefore yields
\(\mu'=iL/(2\pi)-ig_2/(6\pi P)\).
At a zero of \(\wp\), the same heat equation and the eta derivative
give
\[
 \frac{d}{d\tau}\log(\vartheta_1(\mu,\tau)/\eta_D(\tau))
       =L\mu'+\frac{L^2}{4\pi i}.
\]
Differentiating \(D=\operatorname{Im}\beta-
6(\operatorname{Im}\mu)^2/\operatorname{Im}\tau\), and writing
\(b=\operatorname{Im}\mu/\operatorname{Im}\tau\), now gives
\[
 \partial_\tau D=\frac3\pi Z(\mu'-b)
                         -\frac{3i}{4\pi^2}Z^2,
 \qquad \mu'-b=\frac{iZ}{2\pi}-\frac{ig_2}{6\pi P}.
\]
Substitution proves the formula. Where \(dt/d\tau\ne0\), the chain
rule transfers vanishing of this derivative to the coarse base.
The program \texttt{d\_push\_gradient.py} cross-checks the identity
against a separate symmetric difference of the theta-product norm;
that numerical check is not needed for the algebraic proof.
\end{proof}

\begin{proposition}[D.4.5; Proved]
At the equianharmonic point \(\tau=\rho=e^{\pi i/3}\), with the
original period representative \(\mu=(2-\rho)/3\), one has
\[
 \mathcal N(\mu,\rho)=3^{1/4},\qquad
 D_{\rm part}(z_1)=-\frac3{2\pi}\log48.
\]
Consequently the zero locus passes through the coarse point \(t=0\)
exactly at the height
\[
 H_0=\frac3{2\pi}\log48
 =H_*+\frac3{2\pi}\log\frac43.
\]
It has \(D_c(0)<0\) below this height and \(D_c(0)>0\) above it.
In the fixed coarse coordinate the function is \(C^1\) at zero and
\[
 D_{\rm part}(t)=D_{\rm part}(0)
       +\frac7{32\pi}\operatorname{Re}t+O(|t|^{4/3}),\qquad
 \partial_tD_{\rm part}(t)=\frac7{64\pi}+O(|t|^{1/3}).
\]
Thus at \(H=H_0\) the level set is a regular \(C^1\) arc through
zero, tangent to the imaginary axis. This proves the crossing value
and its local regularity requested in D.4, without using an unproved
closed form for \(\beta^{\rm part}(z_1)\).
\end{proposition}
\begin{proof}
For a fixed modulus write \(T=\operatorname{Im}\tau\) and
\[
 u(z)=\log\mathcal N(z,\tau)
 =\log\left|\frac{\vartheta_1(z,\tau)}{\eta_D(\tau)}\right|
       -\frac{\pi(\operatorname{Im}z)^2}{T}.
\]
The theta transformation law makes \(u\) periodic for
\(\Lambda=\mathbb Z+\mathbb Z\tau\). It is locally integrable,
with \(u(z)=\log|z|+c_\tau+o(1)\) at zero, and has no other
singularity on the torus. We first fix its additive normalization
exactly. In the Jacobi product, with
\(Q=e^{2\pi i\tau}\) and \(X=e^{2\pi iz}\), the quotient is
\[
 \frac{\vartheta_1(z,\tau)}{\eta_D(\tau)}
 =iQ^{1/12}X^{-1/2}(1-X)
       \prod_{m\ge1}(1-Q^mX)(1-Q^mX^{-1}),
\]
up to a constant of absolute value one. For
\(z=x+iy\), \(0<y<T\), every product factor has horizontal
mean logarithm zero, by the power series for
\(\log(1-ae^{2\pi ix})\), \(|a|<1\). The series converges
uniformly on closed substrips. Thus
\[
 \int_0^1u(x+iy)\,dx
       =-\frac{\pi T}{6}+\pi y-\frac{\pi y^2}{T}.
\]
The logarithmic singularities at the strip endpoints are
integrable; alternatively one integrates first on closed substrips
and passes to the limit. Integration in \(y\) gives
\(\int_{\mathbb C/\Lambda}u\,dA/T=0\). A tilted fundamental
parallelogram has the same horizontal averages, by period one.
This proves the zero-mean normalization directly.

For \(\tau=\rho\) put \(\alpha=1+\rho\). Since
\(\rho^2=\rho-1\), multiplication by \(\alpha\) on
\((1,\rho)\) has integral matrix
\[
 \begin{pmatrix}1&-1\\1&2\end{pmatrix},\qquad \det=3.
\]
Moreover \(\alpha\mu=1\), \(3\mu=2-\rho\in\Lambda\),
and \(\mu\notin\Lambda\). Its kernel on the torus is therefore
\(K=\{0,\mu,-\mu\}\), and \(|\alpha|^2=3\).
The distributional Euclidean Laplacian is
\[
 \Delta u=2\pi\delta_0-\frac{2\pi}{T}.
\]
Indeed the holomorphic theta factor contributes its simple zero,
and the Gaussian contributes the displayed constant. Pulling back
by the covering \(z\mapsto\alpha z\) shows that
\(u(\alpha z)\) and \(\sum_{v\in K}u(z+v)\) have the same
Laplacian: both have one logarithmic mass at every point of \(K\)
and constant term \(-6\pi/T\). Their difference is consequently
a periodic harmonic function. The logarithmic singularities cancel
locally, so the removable-singularity theorem gives an ordinary
harmonic function on the compact torus. It is constant.
Multiplication by \(\alpha\) and translation preserve normalized
Haar measure. The mean just computed therefore makes that constant
zero, proving the exact distribution identity
\[
 \mathcal N(\alpha z,\rho)
       =\prod_{v\in K}\mathcal N(z+v,\rho).
\]
As \(z\to0\), divide by \(\mathcal N(z,\rho)\). The left
ratio tends to \(|\alpha|=\sqrt3\). Since \(\mathcal N\) is
even, the other two factors tend to \(\mathcal N(\mu,\rho)^2\).
Positivity now gives \(\mathcal N(\mu,\rho)=3^{1/4}\).

Finally the exact R3 normalization is
\[
 D_{\rm part}=\frac6\pi
          \log\frac{\mathcal N(\mu,\tau)}{2\sqrt3}.
\]
Substitution gives \(-3\log48/(2\pi)\). The boundary height is
\(H_*=3\log6/\pi\), so the stated difference is
\(3\log(4/3)/(2\pi)>0\). Addition of \(\operatorname{Im}c\)
proves the crossing height and sign assertions.

We verify regularity in the coarse coordinate rather than infer it
from the ramified modulus. Multiplication by \(\rho^2\) is an
isometry of the torus and preserves the zero-mean Green norm just
used. Also \((\rho^2-1)\mu=\rho-1\in\Lambda\). Hence its
nontrivial rotation fixes \(\mu\), and the real gradient of
\(\log\mathcal N\) vanishes there. In the notation of D.4.4 this
says \(Z(\rho)=0\). Write \(P_0=\wp'_\rho(\mu)\ne0\),
\(E=\pi^2E_2/3\), and \(b=\operatorname{Im}\mu/T\).
The identities of D.4.4 give \(\mu'=b\) at \(\rho\), since
\(g_2(\rho)=0\) and \(L=-2\pi ib\). The heat equation, with
\(\mu\) held fixed, gives
\[
 L_\tau=\frac{-P-2EL}{4\pi i},\qquad L_\mu=-E.
\]
When differentiating \(Z\) along \(\mu(\tau)\), its two terms
involving \(E b\) cancel and the derivative of \(b\) vanishes at
\(\rho\). Thus
\[
 \partial_\tau Z(\rho)=\frac{iP_0}{4\pi},\qquad
 \partial_{\bar\tau}Z(\rho)=0.
\]
Ramanujan's identity and the Weierstrass equation yield
\[
 g_2'(\rho)=-\frac{3i}{\pi}g_3(\rho)
            =\frac{3i}{\pi}P_0^2.
\]
For \(v=\tau-\rho\), it follows that
\[
 3ZP-2g_2=-\frac{21iP_0^2}{4\pi}v+O(|v|^2),\qquad
 \partial_\tau D_{\rm part}
       =\frac{21iP_0^2}{64\pi^4}v^2+O(|v|^3).
\]
Here all remainders are real analytic in \(v,\bar v\).
The exact normalization \(t=j/1728\) gives
\[
 t=-\frac{i g_3(\rho)}{\pi^3}v^3+O(v^4)
   =\frac{iP_0^2}{\pi^3}v^3+O(v^4).
\]
For example this follows immediately from
\(t=E_4^3/(E_4^3-E_6^2)\),
\(E_4'(\rho)=-(2\pi i/3)E_6(\rho)\), and
\(g_3=8\pi^6E_6/27\). Dividing the derivative expansions on
\(v\ne0\) proves
\(\partial_tD_{\rm part}=7/(64\pi)+O(|t|^{1/3})\).
The conjugate derivative has the conjugate expansion. The resulting
real Taylor expansion extends the coarse function as a \(C^1\)
function with nonzero real gradient, giving the asserted local arc
by the implicit function theorem. In particular this is not a
singular level at \(t=0\), despite the vanishing modulus derivative
at the order-three ramification point.

The degree-three lattice matrix, its kernel representatives and all
rational derivative coefficients are checked exactly in
\path{code/d_push_equianharmonic.py}. Independent theta-derivative
and finite-difference evaluations at 90 digits are retained in
\path{code/d_push_equianharmonic_crossing.py} and its output;
they support, but do not supply, the analytic proof.
\end{proof}

\begin{remark}[D.4.6; Known]
This identifier recorded an open step at the earlier checkpoint:
the minimum principle in D.4.2 did not exclude additional maxima,
and the exact crossing in D.4.5 did not classify all levels. That
open status is superseded by D.4.7--D.4.8 below. We retain the
record of the unsuccessful reduction. Lin and Wang (2010,
Theorem~1.3) count critical points of the Green function on a fixed
elliptic curve, where \(Z=0\) is its critical equation. That theorem
does not count zeros along \(\wp_\tau(\mu(\tau))=0\) as the modulus
varies, nor does it address \(3ZP-2g_2\). Our completed argument
therefore supplies a separate nonvanishing certificate for both
factors, with analytic local and cusp estimates.
\end{remark}

\begin{proposition}[D.4.7; Proved]
The function $D_{\rm part}$ has no critical point on the punctured
root neighbourhood $0<|\tau-i|\leq3/20000$, no critical point on the
coarse neighbourhood represented by $|\tau-\rho|\leq1/50$, and no
critical point where $\Im\tau\geq2$. At $\rho$ the two factors
$Z$ and $F=3ZP-2g_2$ of D.4.4 vanish, but the coarse gradient is
nonzero, as in D.4.5. The assertions away from $i$ are certified by
the independently implemented Arb calculation
\path{code/d_root_gradient_certificate.py}; all its intervals and
the 400 local rectangles are retained.
\end{proposition}
\begin{proof}
We first treat $i$ by analytic estimates. Use the holomorphic
functions $G(h,i+x)$ and $H(h,i+x)$ of E.4.2. That proof gives
$|G|<20$, $|H|<0.6$ on the polydisc with radii $1/5$ in $h$ and
$1/10$ in $x$, and gives all its roots over $|x|\leq3/20000$ in
$|h|^2<1/20000$. Put
\[
 a=\frac3{2000},\qquad b=\frac1{800},\qquad
 U=(1-a)^{-2}(1-b)^{-1}-1,\qquad
 V=(1-a)^{-1}(1-b)^{-2}-1.
\]
Differentiate the convergent even power series of $G$. Subtracting
the known coefficients $G_x(0,i)=id$ and
$G_{h^2}(0,i)=-g_2(i)/4$, Cauchy's coefficient inequalities give
\[
 |G_x-id|\leq200U<0.853,\qquad
 \frac{|G_h|}{|h|}\leq2(47.30+500V)<98.611.
\]
The first bound follows by summing
$\sum_{j\geq1,k\geq0}j a^{j-1}b^k$ after its constant term
has been removed. The second uses
$\sum_{j\geq0,k\geq1}k a^j b^{k-1}$ with the same removal.
Since $d>10.03$, $G_x$ is nonzero and the zero curve is locally
$x=x(h)$, with
\[
 |x'(h)|\leq K|h|,\qquad
 K=\frac{2(47.30+500V)}{10.03-200U}<10.745<11.
\]
At $h=0$, E.4.2 forces $x=0$; elsewhere $h$ is thus a valid local
coordinate of this curve. For
$R=H-H(0,i)+\pi h^2/2$, its constant and both linear coefficients
in $(x,h^2)$ vanish. The same sums imply
\[
 |R_h|\leq30V|h|,\qquad |R_x|\leq6U.
\]
Write $T=1+\Im x$ and $y=\Im h$. The norm on the zero curve is
$\log\mathcal N=\Re H-\pi y^2/T$. Its Wirtinger derivative satisfies
\[
 \partial_h\log\mathcal N
 =-\frac\pi2\bar h+\mathcal E,
\]
\[
 \frac{|\mathcal E|}{|h|}
 \leq\frac{30V+6UK}{2}
 +\frac{\pi\epsilon}{1-\epsilon}
 +\frac{\pi K}{40000(1-\epsilon)^2}
 <0.198862<\frac\pi2,\qquad \epsilon=\frac3{20000}.
\]
Indeed the last two terms bound respectively
$\pi i y(T^{-1}-1)$ and
$-\pi i y^2x'(h)/(2T^2)$. This proves strict nonvanishing for
$h\ne0$. Since $D=(6/\pi)\log(\mathcal N/(2\sqrt3))$, it proves
the assertion at $i$ without differentiating a numerical upper bound.

We give the details of the other two certificates. Write $q=e^{2\pi
i\tau}$ and use $\wp,\wp'$ as in E.4.1, together with
\[
\begin{aligned}
 L&=\pi\cot\pi z+
       4\pi\sum_{n\geq1}\frac{q^n}{1-q^n}\sin(2\pi nz),\\
 E&=\frac{\pi^2}{3}-8\pi^2\sum_{n\geq1}\frac{nq^n}{1-q^n},&
 g_2&=\frac{4\pi^4}{3}
       +320\pi^4\sum_{n\geq1}\frac{n^3q^n}{1-q^n}.
\end{aligned}
\]
These follow directly by differentiating the Jacobi product and the
Eisenstein series. The code sums 40 terms. For $0\leq r<1$, set
\[
 V_j(r)=41^j r^{41}M_j(r),\qquad
 (M_0,M_1,M_2,M_3)=
 \left(\frac1{1-r},\frac1{(1-r)^2},
 \frac{1+r}{(1-r)^3},\frac{1+4r+r^2}{(1-r)^4}\right).
\]
Then $\sum_{n>40}n^j r^n\leq V_j(r)$, by
$(40+k)^j\leq41^jk^j$. If $r\geq|q|$ and
$R=r\exp(2\pi|\Im z|)<1$, the tails for
$(\wp,\wp',L,E,g_2)$ are enclosed by complex squares whose
coordinate radii are respectively
\[
 \frac1{1-r}\bigl(
 8\pi^2(V_1(r)+V_1(R)),\ 16\pi^3V_2(R),\
 4\pi V_0(R),\ 8\pi^2V_1(r),\ 320\pi^4V_3(r)\bigr).
\]
They enclose the corresponding complex discs and so are conservative.
Every use of a strict inequality is an Arb comparison between
outward-rounded real balls at 256-bit precision.

For the cusp, let
\[
 z_\infty=\frac12-\frac{i}{2\pi}\log(5+2\sqrt6).
\]
The $q=0$ trigonometric function has a simple zero there. Let $z_0$
be its stored exact dyadic approximation and let $A$ be the stored
dyadic approximation of the inverse derivative at $q=0$. The
closed square $|\Re(z-z_0)|,|\Im(z-z_0)|\leq1/1000$ is invariant
under $z\mapsto z-A\wp(z,q)$ for every $q$ in the square
$|\Re q|,|\Im q|\leq e^{-4\pi}$. The certified supremum-norm
Lipschitz bound is less than $0.024725$ and the image radius is less
than $0.000122623$. Banach's theorem gives the unique analytic root
there, agreeing with the chosen cusp branch. The square in $q$
contains the entire required disc; its larger absolute-value bound
is used in every tail. In the additional interval
$0\leq1/\Im\tau\leq1/2$, direct evaluation gives
\[
 |Z|>1.39894,\qquad |3ZP-2g_2|>127.018.
\]
Consequently neither critical factor vanishes anywhere above height
two. The opposite zero $-z$ changes the sign of $Z$ and $P$ and
leaves $F$ unchanged; translations by periods leave the same test
invariant. Thus the choice of the zero representative does not
restrict this conclusion.

At $\rho$, put $z_\rho=(2-\rho)/3$ and $P_0=\wp'_\rho(z_\rho)$.
The square $\rho+[-1/50,1/50]+i[-1/50,1/50]$ is divided into
$20\times20$ closed rectangles of side $1/500$. This is an exact
rational partition relative to $\rho$. On each rectangle the code
constructs an exact dyadic centre and preconditioner and proves the
same strict contraction criterion for a root square of radius
$1/1000$. Every resulting root lies in the common closed square of
coordinate radius $1/80$ about the dyadic centre of $z_\rho$; the
largest required coordinate radius is less than $0.007731<1/80$.
On this larger square the same uniform derivative calculation
gives $\|1-A\wp'\|_\infty<0.432523<1$ for every modulus in the
whole square about $\rho$. This proves injectivity of $\wp$ there:
if two values agree, the difference of the inputs equals their
difference under $z\mapsto z-A\wp(z)$, contradicting the strict
Lipschitz inequality. This larger square need not be invariant
under that map, and we do not assert that it is. Existence on each
small box and uniqueness in the common square therefore glue one
holomorphic root. The certified nonzero values of $P$ also exclude
a double root. At the centre this is $z_\rho$.

All needed total derivatives on $\wp(\mu)=0$ can be evaluated
without numerical differentiation. With $b=\Im\mu/\Im\tau$,
they are
\[
\begin{aligned}
 \mu'&=\frac{iL}{2\pi}-\frac{i g_2}{6\pi P},\\
 Z_\tau&=\frac{-P-2EL}{4\pi i}-E\mu'
                  +\frac\pi{\Im\tau}(\mu'-b),&
 Z_{\bar\tau}&=\frac\pi{\Im\tau}(b-\overline{\mu'}),\\
 P_\tau&=\frac{-3EP-Lg_2/2}{2\pi i}-\frac{g_2}2\mu',&
 g_2'&=\frac{2iEg_2+3iP^2}{\pi},\\
 F_\tau&=3Z_\tau P+3ZP_\tau-2g_2',&
 F_{\bar\tau}&=3Z_{\bar\tau}P.
\end{aligned}
\]
The first two follow from D.4.4 and the theta heat equation. For the
third, differentiate its displayed partial-$\tau$ identity for
$\wp$ in $\mu$, then add $\wp''\mu'=-g_2\mu'/2$.
The fourth is Ramanujan's identity together with $P^2=-g_3$ at a
zero of $\wp$. In particular these formulas are identities for
the chosen analytic root, not independent-variable derivatives.

Let $A_0=iP_0/(4\pi)$ and $B_0=-21iP_0^2/(4\pi)$. The
400 strict enclosures give, throughout the square,
\[
\begin{aligned}
 |Z_\tau-A_0|+|Z_{\bar\tau}|&<0.423916
                         <2.279897<|A_0|,\\
 |F_\tau-B_0|+|F_{\bar\tau}|&<337.953
                         <1371.701<|B_0|.
\end{aligned}
\]
Since $Z(\rho)=F(\rho)=0$ by D.4.5, integrate each derivative
along the segment from $\rho$ to $\tau$. Its difference from
$A_0(\tau-\rho)$, respectively $B_0(\tau-\rho)$, has strictly
smaller absolute value than that linear term. Hence both factors
are nonzero for $0<|\tau-\rho|\leq1/50$. At the centre itself,
D.4.5 gives the nonzero coarse derivative $7/(64\pi)$.
This completes all local and cusp assertions.
\end{proof}


\begin{proposition}[D.4.8; Proved]
For every \(H>H_*\), the degeneracy locus \(Z_H\) is exactly one
embedded \(C^1\) Jordan curve in the finite \(t\)-plane. It is
real analytic unless \(H=H_0=3\log48/(2\pi)\); at that height it
is real analytic away from \(t=0\) and has the regular crossing
described in D.4.5. There are no singular coarse levels at any
supercritical height. The bounded complementary disc contains
\(t=1\) and has \(D_c>0\); the complementary disc in
\(\mathbb P^1\), containing the cusp, has \(D_c<0\). The point
\(t=0\) is in the latter disc for \(H<H_0\), on the curve for
\(H=H_0\), and in the former disc for \(H>H_0\).

The bounded discs are strictly nested as \(H\) increases, shrink
to \(\{1\}\) as \(H\downarrow H_*\), and exhaust the finite
plane as \(H\to\infty\). Their small-height boundary shape is
the expansion in D.4.1. The maximal open base locus of the complex
manifold construction consists of these two complementary discs.
Approaching their common boundary, the covolume of the marked
root torus lattice is
\[
 \operatorname{covol}(\Pi_c\mathbb Z^4)
       =\operatorname{Im}\tau\,|D_c|\longrightarrow0.
\]
At the boundary the period map has rank three and the exact quotient
and Hausdorff reflection are those of D.4.3. Thus this statement
classifies the locus and the local degeneration requested in D.4.
\end{proposition}
\begin{proof}
We first finish the nonvanishing calculation on the compact part of
the modular domain. Modular invariance of \(\mathcal N\), and
\(\mathcal N(\bar z,-\bar\tau)=\mathcal N(z,\tau)\), reduce it to
\[
 0\leq\Re\tau\leq\tfrac12,\qquad |\tau|\geq1,
 \qquad \Im\tau\leq2.
\]
The reflection identity follows by conjugating the theta and eta
series and the Gaussian. Both roots of \(\wp_\tau=0\) give the
same norm, since they differ by sign and periods. Likewise, a period
translation leaves \(Z,P,F\) unchanged, while replacing \(\mu\)
by \(-\mu\) changes the signs of \(Z,P\) and leaves
\(F=3ZP-2g_2\) unchanged. Thus checking either certified root is
sufficient.

The exact cover audit in E.4 partitions
\([0,1/2]\times[17/20,2]\) into 19,701 root-certified rational
rectangles, 36 omitted rectangles contained in
\(|\tau-i|<3/20000\), and 174 rectangles strictly below
\(|\tau|=1\). It proves equality of the union with that rectangle
and disjointness of interiors by a rational vertical sweep, not by
sampling. The 19,701 E.4 contraction enclosures are the input to
\path{code/d_push_global_arb.py}; no exploratory root approximation
is used as an enclosure. Their stored decimal balls are parsed with
outward error, and the enclosing root radius is increased to its
upper endpoint.

For explicit reproducibility, the code sums the series for
\(L,P,E_2,E_4\) through \(N=36\). Put
\(r=|q|_{\rm upper}\),
\(R=r\exp(2\pi|\Im\mu|_{\rm upper})\), and
\[
 T_j(a)=37^j a^{37}M_j(a),\qquad
 (M_0,M_1,M_2,M_3)=
 \left(\frac1{1-a},\frac1{(1-a)^2},
 \frac{1+a}{(1-a)^3},\frac{1+4a+a^2}{(1-a)^4}\right).
\]
The elementary inequality \((36+k)^j\leq37^jk^j\), for
\(k\geq1\), bounds the corresponding tails. Since
\(|\sin(2\pi n\mu)|\leq e^{2\pi n|\Im\mu|}\) and
\(|1-q^n|\geq1-r\), the absolute errors in \(L,P,E_2,E_4\)
are respectively at most
\[
 \frac1{1-r}\bigl(
 4\pi T_0(R),\ 16\pi^3T_2(R),\
 24T_1(r),\ 240T_3(r)\bigr).
\]
The program verifies \(r,R<1\) on every box and adds these errors
as complex squares. It then forms
\(Z=L+2\pi i\Im\mu/\Im\tau\),
\(g_2=4\pi^4E_4/3\), and \(F=3ZP-2g_2\) using outward-rounded
Arb arithmetic at 192 bits. A box is accepted only when the
certified lower bounds for both absolute values are positive. It
is assigned to a local calculation only when its entire closed
rectangle has distance strictly smaller than the reserved radius
from \(i\) or \(\rho\).

The initial pass accepts 18,263 rectangles, assigns 1,425 to the
\(i\)-disc, and leaves 13 undecided near \(\rho\).
\path{code/d_push_global_refine.py} bisects both rational coordinate
intervals of an undecided rectangle and recomputes its strict root
contraction. Eighteen such subdivision nodes, of depth at most two,
replace those 13 rectangles by 67 leaves: 56 pass and 11 lie wholly
in \(|\tau-\rho|<1/50\). No leaf remains undecided. The resulting
18,319 nonvanishing leaves satisfy the weaker exact rational bounds
\[
 |Z|>\frac{13}{10000},\qquad |F|>13;
\]
the least computed lower endpoints are greater than
\(0.00133357953708\) and \(13.58699524669\), respectively.
All 74 new root contractions, every rational child rectangle, and
every interval output are retained in the refinement JSON file.
\path{code/d_push_global_cover_audit.py} checks the subdivision tree
exactly, checks the cover against E.4, rechecks every local-disc
containment and the displayed rational bounds, and checks the input
hashes. Its output is \texttt{PASS}, with 19,755 final leaves:
18,319 nonvanishing leaves, 1,425 \(i\)-leaves and 11
\(\rho\)-leaves. Together with D.4.7 these are exhaustive.

We now pass from these estimates to the coarse function. Away from
\(t=0,1\), the modulus is a local holomorphic coordinate and
\(P\ne0\): the Weierstrass equation gives \(P^2=-g_3\), and
\(g_3=0\) on a nonsingular elliptic curve gives \(j=1728\).
Consequently D.4.4 and the certificates show that the real
differential of \(D_{\rm part}\) never vanishes there. D.4.7
supplies the same assertion on the punctured root neighbourhood
of \(t=1\). At \(t=0\), D.4.5 supplies a \(C^1\) coarse
extension with nonzero differential. Hence
\[
 f=-D_{\rm part}:\mathbb C\setminus\{1\}
               \longrightarrow(H_*,\infty)
\]
is a \(C^1\) submersion. It is proper: a closed bounded interval
of values excludes a neighbourhood of \(1\), because
\(f(1)=H_*\), and excludes a neighbourhood of infinity, because
\(f\to\infty\) there. Its inverse image is therefore a closed
subset of a compact annulus. The strict minimum at \(1\) is the
independently certified R3 maximum in E.4.

We give the needed level transport also at \(t=0\), without
assuming that this coarse function is \(C^2\). Fix
\(H_*<a<b<\infty\). On the compact band
\(K=f^{-1}([a,b])\), the continuous covector \(df\) is nowhere
zero. At each point choose a constant real direction on which
\(df\) is positive; it remains positive in a neighbourhood.
A finite smooth partition of unity gives a smooth vector field
\(V\) on a neighbourhood of \(K\) with \(df(V)>0\) throughout
\(K\). Compactness gives a positive uniform lower bound. The
smooth flow of \(V\), starting on \(f^{-1}(a)\), therefore reaches
every intermediate level, and then level \(b\), in finite time.
It cannot cease earlier, since before exiting the band it remains
in a compact subset of the vector field's domain. The hitting time
is unique by strict monotonicity of \(f\), and depends \(C^1\) on
the initial point and target height by the implicit function theorem.
Reversing the flow gives the inverse. Thus every two levels in the
band are \(C^1\)-diffeomorphic. This argument applies to arbitrary
finite bands and crosses \(H_0\) without a regularity gap.

For \(a-H_*\) small, D.4.1 proves that \(f^{-1}(a)\) is one
embedded circle. Transport proves that every \(f^{-1}(H)\) is
one embedded circle. The real-analytic implicit function theorem
applies away from \(t=0,1\); the level never contains \(1\),
and it contains \(0\) exactly when \(H=H_0\). This proves the
regularity assertions, including the absence of singular coarse
levels. By the Jordan curve theorem there are exactly two
complementary discs in the sphere. Continuity and the absence of
any other zero imply a constant sign on each disc. The values near
\(1\) and infinity identify their signs. Also
\(D_c(0)=H-H_0\), proving the claimed position of zero. The
inequalities \(f<H\), together with the unique minimum and
properness just proved, give strict nesting, shrinking and
exhaustion.

Finally, in the ordered real coordinates
\((\Re\zeta_1,\Im\zeta_1,\Re\zeta_2,\Im\zeta_2)\) and the
specified lattice basis \((\widehat\gamma,\widehat u,
\widehat w,\widehat\delta)\), direct expansion of the real
period determinant gives
\[
 \det_{\mathbb R}\Pi_c
 =\Im\tau\,\Im\beta_c-6(\Im\mu)^2
 =\Im\tau\,D_c.
\]
The audit program also checks this polynomial identity exactly.
It proves the covolume assertion on either side of the curve. At
every point on it the remaining three independent columns give
rank three, so D.4.3 applies. D.4.2 proves that the total quotient
cannot be a Hausdorff complex manifold at any point of the curve,
including a point where the fixed-fibre period group is discrete.
It also proves that both complementary regions do give complex
manifolds, with the existing Mumford completion on the cusp side.

The literature comparison concerns the exact function just
classified. Chen, Kuo, and Lin (2019, Appendix~A, (A.1) and
(A.3)) give the classical modulus derivatives of \(\wp\) and
\(g_2\); with their \(\eta_1=\pi^2E_2/3\), these agree with our
heat-equation calculation. Their Theorems~1.1--1.4 concern different
holomorphic simple-zero functions and Green-function degeneracy
curves. Chen, Fu, and Lin (2025a, Theorem~1.1; 2025b,
Theorem~1.2) establish generic nondegeneracy for multiple Green
functions and nondegeneracy of the extra critical points of a fixed
torus Green function. None of these statements is the composite
modulus-level assertion above. Searches for the theta norm at
\(\wp=0\), modulus critical points, and Green-function degeneracy
curves, including those primary papers and their references, did
not locate this precise global classification. We claim no priority;
the proof here consists of the stated analytic identities, finite
interval certificate, and level-transport argument.
\nocite{Buchdahl1999,Grauert60push,Grauert63push,CornulierChabauty,LinWang2010D,ChenKuoLin2019D,ChenFuLin2025GenericD,ChenFuLin2025HitchinD}
\end{proof}
